\documentclass[10pt,a4paper]{amsart}
\usepackage[margin=19mm]{geometry}
\usepackage{amsmath,amssymb,mathtools}
\usepackage[scr=rsfs,cal=euler]{mathalfa}
\usepackage{xfrac}
\usepackage[unicode,hidelinks,bookmarks=false]{hyperref}
\newcommand{\ie}{i.e.\ }
\newcommand{\cf}{cf.\ }
\newcommand{\resp}{respectively\ }
\newcommand{\Subsection}{\subsection}
\providecommand{\nobreakdash}{\nobreak\hbox{-}}
\setlength{\emergencystretch}{2em}
\multlinegap0pt
\usepackage{amssymb}
\usepackage{algorithm}\usepackage{algpseudocode}
\usepackage{xcolor}
\usepackage{enumerate}
\usepackage{bm}
\usepackage{dsfont}
\usepackage{mathtools}
\newtheorem{theorem}{Theorem}
\newtheorem{proposition}{Proposition}
\DeclareMathOperator{\Gal}{Gal}
\newcommand{\Z}{\mathbb{Z}} %integers
\newcommand{\de}{{\mathbf{n}}}  %temporary symbol for the degree of $f$
\newcommand{\os}{{\mathbf{m}}}  %temporary symbol for the order of $\sigma$
\title{The first tight classification of skew-constacyclic codes over finite fields}
\author{Monica Nevins, Susanne Pumluen}
\date{Source: arXiv:2608.21339v1 (CC BY 4.0)}
\begin{document}
\maketitle

\noindent\emph{Source and licence.} The first tight classification of skew-constacyclic codes over finite fields. Version arXiv:2608.21339v1 (CC BY 4.0), distributed by arXiv under the Creative Commons Attribution 4.0 International licence, http://creativecommons.org/licenses/by/4.0/, as recorded in the arXiv licence metadata for this exact version and shown on the abstract page https://arxiv.org/abs/2608.21339v1. Full licence notice: \texttt{licenses/arxiv-2608.21339-CC-BY-4.0.txt}.\par\smallskip

\noindent\textit{Editorial scope.} Counted unit: the article's proposition T:F-tildeequivalencecount (source0.tex lines 661--666). The article's complete original proof of it is delivered with it (source0.tex lines 668--693, one \texttt{proof} environment of 26 lines). No mathematics is written, repaired or re-derived: the statement and the proof are the article's own lines, in the article's own notation, and no retained line is substituted or re-flowed.  Retained uncounted as the source-local context the proof needs: lines 564--566.  Nothing was omitted from any retained span, so the package carries no omission record.  No retained line contains a citation, so the wrapper carries no bibliography and no bibliography entry.  No retained line contains a figure environment, an image-inclusion command or drawing code, so no image is delivered and the package declares no asset dependency.  Editorial formatting only, in addition: an AMS \texttt{amsart} preamble that restates the article's own declarations and macros used by the retained lines, namely the environments \texttt{theorem}, \texttt{proposition} on separate counters, together with the standard packages those lines need.  The retained environments print in this document as Theorem 1, Proposition 1 in order of appearance, whereas the article prints the same items with the numbering of its own full document.  The complete native source of the article as distributed in the arXiv e-print for this exact version is bundled unchanged as \texttt{sources/208242/source0.tex}.
\begin{theorem}\label{L:Chenclassesfinitefield}
    The distinct Chen-equivalence classes of  ambient Petit rings of the form  $K[t;\sigma]/K[t;\sigma](t^\de-a)$, where $[K:F]=\os$, are parametrized by taking $a\in \{1, \xi, \cdots, \xi^{s-1}\}$, where $s = \gcd\left(q^\os-1, \frac{q^\de-1}{q-1}\right).$
    \end{theorem}

\begin{proposition}\label{T:F-tildeequivalencecount}
    Let $\tilde{F}$ be a subfield of $F$ and denote its size by $\tilde{q}$.  The number of equivalence classes over $\tilde{F}$ of algebras of the form $K[t;\sigma]/K[t;\sigma](t^\de-a)$, as $a$ ranges over $K^\times$ is
    $$
    \sum_{d|s}\frac{\varphi(d)}{o(\tilde{q})_d}.
    $$
\end{proposition}

\begin{proof}
    Let $N=[K:\tilde{F}]=\log_p(q^\os)/\log_p(\tilde{q})$ denote the degree of this extension.  The automorphism of $K$ defined by $\tau(x)=x^{\tilde{q}}$ for each $x\in K$ generates the cyclic  Galois group ${\rm Gal}(K/\tilde{F})$ and commutes with $\sigma$.  In particular, $\tau$ stabilizes $K^\times$ and the subgroup $N^\sigma_\de(K^\times)$, so acts as a permutation on  the cosets of $N^\sigma_\de(K^\times)=\langle \xi^s\rangle$, where $s=\gcd\left(q^\os-1, \frac{q^\de-1}{q-1}\right)$.
    By Lemma~\ref{L:Chenclassesfinitefield}, each coset is represented by $\xi^i$ for some $0\leq i <s$.  For each such $i$, the orbit of $\xi^i$ under the Galois group is
    $$
    \Gal(K/\tilde{F})\cdot \xi^i = \{ \xi^{i\tilde{q}^r} \mid 0\leq r < N\}.
    $$
    Thus the Galois orbit of the coset $\xi^i N^\sigma_\de(K^\times) = \xi^i\langle \xi^s\rangle$ in $K^\times$ is the union
    $$
    \Xi_i:=\bigcup_{0\leq r<N}\xi^{i\tilde{q}^r}N^\sigma_\de(K^\times).
    $$
    These orbits need not be distinct: for example if $j\equiv i\tilde{q}\mod s$ then $\xi^j$ and $\xi^{i\tilde{q}}$ represent the same coset, so $\Xi_i=\Xi_{j}$.  The orbits are also not all of the same size: for example, the trivial coset forms its own orbit, whereas in general the size depends on the order of $\tilde{q}$ mod $s$.  Nevertheless, each of the \emph{distinct} orbits  $\Xi_i$ defines a single equivalence class of algebras $K[t;\sigma]/K[t;\sigma](t^\de-a)$ over $\tilde{F}$, namely, by letting $a$ run over $\Xi_i$.

    To count the number of distinct orbits, we proceed as follows.  Let $d$ be a divisor of $s$.  Set ${\mathbf R}_d = \{i\in \Z/s\Z\mid \gcd(i,s)=s/d\}$ to be the set of all elements of additive order exactly $d$; since $\gcd(\tilde{q},s)=1$, we have that for each $i\in {\mathbf R}_d$ and $0\leq r<N$ that $\gcd(i\tilde{q}^r,s)=s/d$ so
    $$
    \Xi(d) = \bigcup_{j\in \mathbf{R}_d} \xi^j N^\sigma_\de(K^\times)
    $$
    is the $\Gal(K/\tilde{F})$-invariant union of the orbits $\Xi_i$ for which $\gcd(i,s)=s/d$.  Moreover, since the sets $\mathbf{R}_d$ partition $
    \Z/s\Z$, we have
    $$
    K^\times = \bigsqcup_{d|s} \Xi(d)
    $$
    so it suffices to count the number of distinct $\Gal(K/\tilde{F})$ orbits in each such $\Xi(d)$.


    Fix $d$ dividing $s$.  Since  $\gcd(d,\tilde{q})=1$, we may set $s_0:=o(\tilde{q})_d$ to be the order of $\tilde{q}$ in $(\Z/d\Z)^\times$. Thus for each $i\in {\mathbf R}_d$,  $s_0$ is the least power of $\tilde{q}$ for which we have $i(\tilde{q}^{s_0}-1)\equiv 0\mod s$, whence it is the least power such that $i\tilde{q}^{s_0}\equiv i\mod s$. Therefore for each $i\in {\mathbf R}_d$, the number of distinct cosets in the orbit $\Xi_i$ is $s_0$.  Since the map $k\mapsto k(s/d)$ is an isomorphism of $(\Z/d\Z)^\times$ onto ${\mathbf R}_d$, $\Xi(d)$ contains precisely $\varphi(d)$ distinct cosets, which are partitioned into $\varphi(d)/s_0$ orbits of size $s_0=o(\tilde{q})_d$, as required.
\end{proof}

\end{document}
